\documentclass[11pt]{article}
\usepackage[margin=1in]{geometry}
\usepackage{amsmath,amssymb,amsthm,mathtools}
\usepackage[T1]{fontenc}
\usepackage[utf8]{inputenc}
\usepackage{enumitem}
\usepackage{booktabs,longtable,array}
\usepackage{hyperref}
\usepackage{xcolor}
\usepackage{microtype}
\hypersetup{colorlinks=true, linkcolor=blue!50!black, citecolor=blue!50!black, urlcolor=blue!50!black}

\newtheorem{theorem}{Theorem}[section]
\newtheorem{proposition}[theorem]{Proposition}
\newtheorem{lemma}[theorem]{Lemma}
\newtheorem{corollary}[theorem]{Corollary}
\newtheorem{claim}[theorem]{Claim}
\theoremstyle{definition}
\newtheorem{definition}[theorem]{Definition}
\newtheorem{assumption}[theorem]{Assumption}
\theoremstyle{remark}
\newtheorem{remark}[theorem]{Remark}

\newcommand{\KL}{\mathrm{KL}}
\newcommand{\Ent}{\mathrm{Ent}}
\newcommand{\Prob}{\mathsf{P}}
\newcommand{\E}{\mathbb{E}}
\newcommand{\R}{\mathbb{R}}
\newcommand{\N}{\mathbb{N}}
\newcommand{\Pcal}{\mathcal{P}}
\newcommand{\Mcal}{\mathcal{M}}
\newcommand{\Qcal}{\mathcal{Q}}
\newcommand{\Zcal}{\mathcal{Z}}
\newcommand{\Fcal}{\mathcal{F}}
\newcommand{\Gcal}{\mathcal{G}}
\newcommand{\W}{\mathcal{W}}
\newcommand{\one}{\mathbf{1}}
\newcommand{\dd}{\,d}
\newcommand{\supp}{\operatorname{supp}}
\newcommand{\argmin}{\operatorname*{arg\,min}}
\newcommand{\argmax}{\operatorname*{arg\,max}}
\newcommand{\esssup}{\operatorname*{ess\,sup}}
\newcommand{\ri}{\operatorname{ri}}
\newcommand{\dom}{\operatorname{dom}}
\newcommand{\Law}{\operatorname{Law}}
\newcommand{\Var}{\operatorname{Var}}
\newcommand{\Cov}{\operatorname{Cov}}
\newcommand{\diag}{\operatorname{diag}}
\newcommand{\Id}{\mathrm{Id}}
\newcommand{\Lip}{\operatorname{Lip}}
\newcommand{\osc}{\operatorname{osc}}
\newcommand{\diam}{\operatorname{diam}}
\newcommand{\proj}{\operatorname{proj}}
\newcommand{\e}{\mathrm{e}}

\newenvironment{argument}{\paragraph{Argument.}\begin{enumerate}[leftmargin=2em]}{\end{enumerate}}
\newenvironment{proofoutline}{\paragraph{Proof architecture.}\begin{enumerate}[leftmargin=2em]}{\end{enumerate}}


% Source-faithfulness apparatus used by the paper reconstructions.
\newcommand{\sourceanchor}[1]{\par\smallskip\noindent\textbf{Source anchor.} #1\par\smallskip}
\newcommand{\sourcecomment}[1]{\begin{remark}#1\end{remark}}
\newenvironment{sourceguide}
  {\begin{description}[style=nextline,leftmargin=3.2cm,labelwidth=3cm]}
  {\end{description}}

% Backward-compatible environments used by some of the older reconstructions.
\newenvironment{distilled}[1][]{\begin{theorem}[#1]}{\end{theorem}}
\newenvironment{distillation}{\begin{enumerate}[leftmargin=2em]}{\end{enumerate}}

\newenvironment{commentarynotes}{\begin{remark}\begin{enumerate}[leftmargin=2em]}{\end{enumerate}\end{remark}}

\title{Modernized Proof Notes for Yeh (1978): Conditional Wiener Integrals under Translation}
\author{Generated proof-target note for the AIQ DRSB formalization corpus}
\date{Source PDF: \texttt{Transformations of Wiener Integrals Under Translations-Yeh-1977.pdf}}
\begin{document}
\maketitle

\section*{Reference and purpose}
This note expands J. Yeh, ``Transformation of Conditional Wiener Integrals under Translation and the Cameron--Martin Translation Theorem,'' \emph{Tohoku Mathematical Journal}, 30(4):505--515, 1978 \cite{Yeh1978}.  This is a non-arXiv source.  It is useful for Brownian bridge and endpoint-conditioned path-space identities.

\section{Wiener space and conditional Wiener integrals}
Fix $t>0$.  Let
\[
  C[0,t]=\{x:[0,t]\to\R: x\hbox{ continuous},\ x(0)=0\}
\]
with Wiener measure $m_w$.  Cylinder sets are generated by evaluations
\[
  x\mapsto (x(s_1),\ldots,x(s_n)),
  \qquad 0=s_0<s_1<\cdots<s_n\le t,
\]
with the usual Gaussian density
\[
  (2\pi)^{-n/2}\prod_{j=1}^n(s_j-s_{j-1})^{-1/2}
  \exp\left(-\frac12\sum_{j=1}^n\frac{(\xi_j-\xi_{j-1})^2}{s_j-s_{j-1}}\right).
\]
Yeh conditions on the endpoint functional
\begin{equation}
  X[x]=x(t). \tag{X}\label{eq:yeh-X}
\end{equation}
The law of $X$ is $N(0,t)$ on $(\R,\mathcal B(\R))$:
\begin{equation}
  P_X(B)=(2\pi t)^{-1/2}\int_B e^{-\xi^2/(2t)}\,d\xi. \tag{PX}\label{eq:yeh-PX}
\end{equation}
For an integrable functional $Y$ on Wiener space, a conditional Wiener integral $E^w(Y\mid X)$ is a version of the conditional expectation as a function of the endpoint $\xi\in\R$, i.e.
\[
  \int_{X^{-1}(B)}Y(x)m_w(dx)=\int_BE^w(Y\mid X)(\xi)P_X(d\xi).
\]

\section{An abstract conditional-expectation transformation theorem}
Before specializing to Wiener space, Yeh proves a general theorem about conditional expectation under a measurable transformation.
Let $(\Omega,\mathcal B,P)$ be a probability space, $X:\Omega\to S$ a measurable map, and $T:\Omega\to\Omega$ a measurable transformation.

\subsection{Lemma 1: changing the conditioning map}
Assume there is a measurable bijection $h:S\to S$ with measurable inverse such that
\begin{equation}
  X\circ T=h\circ X \quad\hbox{a.e.} \tag{commute}\label{eq:yeh-commute}
\end{equation}
Then for any integrable $Z$,
\begin{equation}
  E(Z\mid X\circ T)(\xi)=E(Z\mid X)(h^{-1}\xi)
  \quad\hbox{for }P_{X\circ T}\hbox{-a.e. }\xi. \tag{L1}\label{eq:yeh-L1}
\end{equation}
The proof is purely measure-theoretic.  For $F\subset S$ measurable,
\[(X\circ T)^{-1}(F)\]
and $(h\circ X)^{-1}(F)=X^{-1}(h^{-1}F)$ differ by a null set.  The defining integral identity for conditional expectation then identifies the two versions after pushforward by $h$.

\subsection{Lemma 2: change of variables by a Radon--Nikodym density}
Let $(\Omega,\mathcal B,\mu)$ be finite and $T:\Omega\to\Omega$ measurable.  If there exists a measurable $J$ such that
\begin{equation}
  \mu(A)=\int_{T^{-1}A}J(\omega)\mu(d\omega) \quad(A\in\mathcal B), \tag{RN}\label{eq:yeh-RN}
\end{equation}
then for every measurable $Z$,
\begin{equation}
  \int_\Omega Z(\omega)\mu(d\omega)
  =\int_\Omega Z(T\omega)J(\omega)\mu(d\omega), \tag{CV}\label{eq:yeh-CV}
\end{equation}
with the usual convention that existence of one side implies existence and equality of the other.
Yeh proves this first for simple functions, then for nonnegative functions by monotone convergence, then for signed functions by positive/negative parts.

\subsection{Theorem 1}
Assume:
\begin{enumerate}[leftmargin=2em]
\item the density/change-of-variables condition \eqref{eq:yeh-RN};
\item the endpoint commutation condition \eqref{eq:yeh-commute};
\item optionally, $P_{X\circ T}\ll P_X$.
\end{enumerate}
Then for measurable $g:S\to\R$,
\begin{equation}
  \int_S g(\xi)E(Z\mid X)(\xi)P_X(d\xi)
  =\int_S g(\xi)E((Z\circ T)J\mid X)(h^{-1}\xi)P_{X\circ T}(d\xi). \tag{T1-int}\label{eq:yeh-T1-int}
\end{equation}
If $P_{X\circ T}\ll P_X$, then arbitrary versions satisfy
\begin{equation}
  E(Z\mid X)(\xi)=E((Z\circ T)J\mid X)(h^{-1}\xi)\frac{dP_{X\circ T}}{dP_X}(\xi)
  \quad P_X\hbox{-a.e.} \tag{T1}\label{eq:yeh-T1}
\end{equation}
This theorem is the abstract machine used in the Wiener translation theorem.

\section{Cameron--Martin translation input}
Let $x_0\in C[0,t]$ be absolutely continuous with
\[
  \int_0^t |x_0'(s)|^2ds<\infty.
\]
Define
\begin{equation}
  T[x]=x+x_0. \tag{trans}\label{eq:yeh-trans}
\end{equation}
The Cameron--Martin theorem gives
\begin{equation}
  m_w(\Gamma)=\int_{T^{-1}\Gamma}J(x)m_w(dx), \tag{CM}\label{eq:yeh-CM}
\end{equation}
where
\begin{equation}
  J(x)=\exp\left(-\frac12\int_0^t|x_0'(s)|^2ds\right)
       \exp\left(-\int_0^t x_0'(s)\,dx(s)\right). \tag{J}\label{eq:yeh-J}
\end{equation}
The second integral is a Paley--Wiener--Zygmund--Riemann--Stieltjes integral in the general $L^2$ derivative case; it agrees with the ordinary Riemann--Stieltjes integral when $x_0'$ has bounded variation.

For the endpoint map $X[x]=x(t)$,
\[
  (X\circ T)[x]=x(t)+x_0(t)=h(X[x]),
  \qquad h(\xi)=\xi+x_0(t),
\]
so
\[
  h^{-1}(\xi)=\xi-x_0(t).
\]
Also
\[
  P_X=N(0,t),
  \qquad
  P_{X\circ T}=N(x_0(t),t),
\]
and therefore
\begin{equation}
  \frac{dP_{X\circ T}}{dP_X}(\xi)
  =\exp\left(-\frac{x_0(t)^2}{2t}\right)
   \exp\left(\frac{\xi x_0(t)}{t}\right). \tag{endpointRN}\label{eq:yeh-endpointRN}
\end{equation}

\section{Theorem 2: transformation of conditional Wiener integrals}
Substituting \eqref{eq:yeh-J} and \eqref{eq:yeh-endpointRN} into the abstract theorem gives Yeh's main theorem.

\begin{distilled}[Yeh Theorem 2]
Let $Y$ be an integrable functional on Wiener space and let $X[x]=x(t)$.  If $x_0$ is absolutely continuous with $x_0'\in L^2[0,t]$, then for arbitrary versions,
\begin{align}
  E^w(Y\mid X)(\xi)
  &=E^w\big(Y[\cdot+x_0]J\mid X\big)(\xi-x_0(t)) 
    \exp\left(-\frac{x_0(t)^2}{2t}\right)
    \exp\left(\frac{\xi x_0(t)}{t}\right) \tag{Y2}\label{eq:yeh-Y2}
\end{align}
for Lebesgue-a.e. $\xi$, where $J$ is the Cameron--Martin density \eqref{eq:yeh-J}.
\end{distilled}

The proof is short after Theorem 1:
\begin{enumerate}[leftmargin=2em]
\item use Cameron--Martin to verify condition 1 of Theorem 1;
\item use $X(Tx)=X(x)+x_0(t)$ to verify condition 2;
\item compute the endpoint Gaussian density ratio \eqref{eq:yeh-endpointRN} for condition 3;
\item apply \eqref{eq:yeh-T1}.
\end{enumerate}

\section{Corollary: conditional exponential integral}
Set $Y\equiv1$ in Theorem 2.  Since $E^w(1\mid X)=1$, the theorem can be solved for the conditional expectation of the exponential factor:
\begin{equation}
  E^w\left(\exp\left[-\int_0^t x_0'(s)\,dx(s)\right]\middle|X\right)(\xi)
  =\exp\left(\frac12\int_0^t|x_0'(s)|^2ds\right)
   \exp\left(-\frac{x_0(t)^2}{2t}\right)
   \exp\left(-\frac{\xi x_0(t)}{t}\right). \tag{cor}\label{eq:yeh-cor}
\end{equation}
This formula is an endpoint-conditioned Cameron--Martin identity.  It is exactly the kind of calculation that appears when rewriting Brownian bridge expectations after a deterministic translation.

\section{Example: exponential of a Paley--Wiener integral}
Yeh considers
\[
  Y[x]=\exp\left(\lambda\int_0^t p(s)\,dx(s)\right),
\]
where $p$ is continuous and of bounded variation.  Choose
\[
  x_0(s)=-\lambda\int_0^s p(r)\,dr,
  \qquad x_0'(s)=-\lambda p(s).
\]
Then \eqref{eq:yeh-cor} gives a version of $E^w(Y\mid X)(\xi)$:
\begin{align}
  E^w(Y\mid X)(\xi)
  &=\exp\left(\frac{\lambda^2}{2}\int_0^t p(s)^2ds\right)
    \exp\left(-\frac{\lambda^2}{2t}\left(\int_0^t p(s)ds\right)^2\right)
    \exp\left(\frac{\lambda\xi}{t}\int_0^t p(s)ds\right). \tag{ex}\label{eq:yeh-ex}
\end{align}
This is the Gaussian bridge regression formula in conditional-Wiener-integral notation.

\section{Formalization relevance}
\begin{formalizationnotes}
\item This paper is relevant only for continuum endpoint-conditioned path measures, especially Brownian bridge translation identities.
\item The abstract Theorem 1 is ordinary measure theory: conditional expectation, pushforward measures, Radon--Nikodym derivative, and a measurable bijection commuting with the conditioning map.  This part is much more plausible to formalize than the Wiener-specific Cameron--Martin theorem.
\item The Wiener-specific theorem should be treated as a theorem seam unless the project develops Gaussian/Wiener measure infrastructure.
\item Formula \eqref{eq:yeh-endpointRN} is finite-dimensional Gaussian algebra and is a good formalization target if endpoint shifts appear in bridge calculations.
\end{formalizationnotes}

\section*{Checklist of details retained}
\begin{itemize}[leftmargin=2em]
\item Definition of conditional Wiener integral as a version over endpoint law $P_X$.
\item Lemma 1: commutation of conditioning map with a measurable bijection.
\item Lemma 2: change of variables from a density relation.
\item Theorem 1: abstract transformation of conditional expectations.
\item Cameron--Martin density used by Yeh.
\item Endpoint Gaussian Radon--Nikodym derivative.
\item Theorem 2 formula with the shifted endpoint argument $\xi-x_0(t)$.
\item Corollary and example computing a conditional exponential integral.
\end{itemize}

\begin{thebibliography}{99}
\bibitem{ChenGeorgiouPavon2021} Yongxin Chen, Tryphon T. Georgiou, and Michele Pavon. \emph{Stochastic Control Liaisons: Richard Sinkhorn Meets Gaspard Monge on a Schr\"odinger Bridge}. SIAM Review, 63(2):249--313, 2021. doi:10.1137/20M1339982. arXiv:2005.10963.
\bibitem{Leonard2014Survey} Christian L\'eonard. \emph{A Survey of the Schr\"odinger Problem and Some of Its Connections with Optimal Transport}. Discrete and Continuous Dynamical Systems--A, 34(4):1533--1574, 2014. doi:10.3934/dcds.2014.34.1533. arXiv:1308.0215.
\bibitem{Fortet1940} Robert Fortet. \emph{Resolution d'un systeme d'equations de M. Schrodinger}. Journal de mathematiques pures et appliquees, 9e serie, 19(1--4):83--105, 1940.
\bibitem{SinkhornKnopp1967} Richard Sinkhorn and Paul Knopp. \emph{Concerning Nonnegative Matrices and Doubly Stochastic Matrices}. Pacific Journal of Mathematics, 21(2):343--348, 1967.
\bibitem{FranklinLorenz1989} Joel Franklin and Jens Lorenz. \emph{On the Scaling of Multidimensional Matrices}. Linear Algebra and its Applications, 114/115:717--735, 1989. doi:10.1016/0024-3795(89)90490-4.
\bibitem{Carroll2004} Joseph E. Carroll. \emph{Birkhoff's Contraction Coefficient}. Linear Algebra and its Applications, 389:227--234, 2004. doi:10.1016/j.laa.2004.02.039.
\bibitem{EvesonNussbaum1995} Simon P. Eveson and Roger D. Nussbaum. \emph{An Elementary Proof of the Birkhoff--Hopf Theorem}. Mathematical Proceedings of the Cambridge Philosophical Society, 117(1):31--55, 1995.
\bibitem{BlanchetMurthy2019} Jose Blanchet and Karthyek R. A. Murthy. \emph{Quantifying Distributional Model Risk via Optimal Transport}. Mathematics of Operations Research, 44(2):565--600, 2019. doi:10.1287/moor.2018.0936. arXiv:1604.01446.
\bibitem{GaoKleywegt2023} Rui Gao and Anton J. Kleywegt. \emph{Distributionally Robust Stochastic Optimization with Wasserstein Distance}. Mathematics of Operations Research, 48(2):603--655, 2023. doi:10.1287/moor.2022.1275. arXiv:1604.02199.
\bibitem{MohajerinEsfahaniKuhn2018} Peyman Mohajerin Esfahani and Daniel Kuhn. \emph{Data-Driven Distributionally Robust Optimization Using the Wasserstein Metric: Performance Guarantees and Tractable Reformulations}. Mathematical Programming, 171(1--2):115--166, 2018. doi:10.1007/s10107-017-1172-1. arXiv:1505.05116.
\bibitem{WangGaoXie2025} Jie Wang, Rui Gao, and Yao Xie. \emph{Sinkhorn Distributionally Robust Optimization}. Operations Research, 2025. doi:10.1287/opre.2023.0294. arXiv:2109.11926.
\bibitem{Girsanov1960} I. V. Girsanov. \emph{On Transforming a Certain Class of Stochastic Processes by Absolutely Continuous Substitution of Measures}. Theory of Probability and Its Applications, 5(3):285--301, 1960. doi:10.1137/1105027.
\bibitem{CameronMartin1945} R. H. Cameron and W. T. Martin. \emph{Transformations of Wiener Integrals under a General Class of Linear Transformations}. Transactions of the American Mathematical Society, 58(2):184--219, 1945.
\bibitem{Yeh1978} J. Yeh. \emph{Transformation of Conditional Wiener Integrals under Translation and the Cameron--Martin Translation Theorem}. Tohoku Mathematical Journal, 30(4):505--515, 1978.
\bibitem{Kakutani1948} Shizuo Kakutani. \emph{On Equivalence of Infinite Product Measures}. Annals of Mathematics, 49(1):214--224, 1948.
\bibitem{ShiDeBortoliCampbellDoucet2023} Yuyang Shi, Valentin De Bortoli, Andrew Campbell, and Arnaud Doucet. \emph{Diffusion Schr\"odinger Bridge Matching}. Advances in Neural Information Processing Systems, 2023. arXiv:2303.16852.
\bibitem{LiuLipmanNickelKarrerTheodorouChen2024} Guan-Horng Liu, Yaron Lipman, Maximilian Nickel, Brian Karrer, Evangelos A. Theodorou, and Ricky T. Q. Chen. \emph{Generalized Schr\"odinger Bridge Matching}. International Conference on Learning Representations, 2024. arXiv:2310.02233.
\bibitem{DomingoEnrichDrozdzalKarrerChen2025} Carles Domingo-Enrich, Michal Drozdzal, Brian Karrer, and Ricky T. Q. Chen. \emph{Adjoint Matching: Fine-Tuning Flow and Diffusion Generative Models with Memoryless Stochastic Optimal Control}. arXiv preprint arXiv:2409.08861, 2025.

\bibitem{Birkhoff1957} Garrett Birkhoff. \emph{Extensions of Jentzsch's Theorem}. Transactions of the American Mathematical Society, 85(1):219--227, 1957.
\bibitem{Sinkhorn1964} Richard Sinkhorn. \emph{A Relationship Between Arbitrary Positive Matrices and Doubly Stochastic Matrices}. Annals of Mathematical Statistics, 35(2):876--879, 1964.
\bibitem{Sinkhorn1967Prescribed} Richard Sinkhorn. \emph{Diagonal Equivalence to Matrices with Prescribed Row and Column Sums}. American Mathematical Monthly, 74(4):402--405, 1967.
\bibitem{Bauer1965} Friedrich L. Bauer. \emph{An Elementary Proof of the Hopf Inequality for Positive Operators}. Numerische Mathematik, 7:331--337, 1965.
\bibitem{Bushell1973} P. J. Bushell. \emph{Hilbert's Metric and Positive Contraction Mappings in a Banach Space}. Archive for Rational Mechanics and Analysis, 52:330--338, 1973.

\end{thebibliography}

\end{document}
